\documentclass[a4paper,11pt]{article}
\usepackage[utf8]{inputenc}
\usepackage[T1]{fontenc}
\usepackage[french]{babel}
\usepackage[left=2cm,right=2cm,top=2cm,bottom=2cm]{geometry}
\usepackage{xcolor}
\usepackage{hyperref}
\usepackage{fontawesome5}
\usepackage{parskip}

% --- Couleurs ---
\definecolor{primary}{RGB}{35, 55, 123}
\hypersetup{colorlinks=true, urlcolor=primary}

\begin{document}

% --- EXPÉDITEUR ---
\noindent
\begin{minipage}[t]{0.5\textwidth}
    \textbf{Loïc Aron MBASSI EWOLO}\\
    Élève-Ingénieur Bac+5 -- Double diplôme ENSEA\\[3pt]
    \faMapMarker\ Cergy, France (Mobile IDF)\\
    \faPhone\ (+33) 7 59 52 33 98\\
    \faEnvelope\ \href{mailto:loicaron.mbassiewolo@ensea.fr}{loicaron.mbassiewolo@ensea.fr}\\
    \faLinkedin\ \href{https://www.linkedin.com/in/loic-mbassi/}{linkedin.com/in/loic-mbassi}
\end{minipage}
\hfill
\begin{minipage}[t]{0.4\textwidth}
    \raggedleft
    Cergy, le 12 mars 2026
\end{minipage}

\vspace{1.2cm}

% --- DESTINATAIRE ---
\hfill
\begin{minipage}[t]{0.5\textwidth}
    \raggedleft
    \textbf{ValoTec}\\
    Service Recrutement / R\&D\\
    Villejuif Bio Park\\
    1 mail du Professeur Georges Mathé\\
    94800 Villejuif
\end{minipage}

\vspace{1cm}

\noindent\textbf{Objet :} Candidature au poste d'Ingénieur junior R\&D et Industrialisation Électronique -- Contrat de Professionnalisation (Septembre 2026)

\vspace{0.8cm}

Madame, Monsieur,

Souder des capteurs IMU et flex sur PCB, valider des signaux à l'oscilloscope, et mettre au point un prototype embarqué -- c'est précisément ce que j'ai fait au laboratoire IOTA de l'École Polytechnique de Yaoundé, sur le projet Harmony Gloves. Cette expérience concrète, combinée à ma formation à l'\textbf{ENSEA} (double diplôme Systèmes Embarqués, VHDL, traitement du signal), me donne envie de contribuer aux projets R\&D de \textbf{ValoTec} autour des dispositifs médicaux actifs.

Dès Septembre 2026, libéré de mon stage de recherche à \textbf{INRIA Saclay} (Équipe TRIBE, plateau de Saclay -- Institut Polytechnique de Paris), je suis disponible pour un \textbf{Contrat de Professionnalisation} correspondant à ma 3\textsuperscript{ème} année à l'ENSEA. Ce stage INRIA, sur l'analyse automatisée des applications mobiles (NLP, Python, rédaction scientifique en anglais), renforce ma capacité à documenter rigoureusement mes travaux et à travailler dans un environnement technique exigeant.

Sur le plan électronique, mon projet le plus représentatif est le développement de gants instrumentés pour la traduction de la langue des signes. J'ai réalisé l'intégration physique des composants, la validation des acquisitions via protocoles I2C/SPI sur STM32, et le déploiement d'un modèle TinyML embarqué. En parallèle, le challenge de robotique autonome CoHoMa (armée française, ENSEA) m'a conduit à travailler en C/C++ sur Nvidia Jetson, avec fusion de données Lidar 3D et caméra de profondeur. Ces projets m'ont appris à progresser rapidement en équipe et à rendre compte de mes travaux par écrit, sous supervision limitée.

La pluridisciplinarité de ValoTec -- électronique, logiciel embarqué, mécanique, industrialisation -- et le travail en binôme avec revue systématique sont des conditions dans lesquelles je progresse efficacement. Je suis disponible dès \textbf{Septembre 2026} pour un Contrat de Professionnalisation, avec l'ambition de m'investir sur la durée et d'évoluer vers un CDI à l'issue.

Je reste à votre disposition pour un entretien afin de vous présenter mes réalisations et de discuter de vos projets en cours.

Je vous prie d'agréer, Madame, Monsieur, l'expression de mes salutations distinguées.

\vspace{0.8cm}

\textbf{Loïc Aron MBASSI EWOLO}\\
\textit{Élève-Ingénieur ENSEA / Polytechnique Yaoundé}

\end{document}
